\chapter{Reference Frames, Light, and the Geometry of Spacetime}

\section{The Course Begins Again with Relativity}

This quarter returns to classical physics, but not to the Newtonian picture in which space and time are independent and universal. The immediate subjects are special relativity and classical field theory, especially electromagnetism. Quantum mechanics will not be used. The earlier mechanics course supplied the action principle and the Hamiltonian and Lagrangian descriptions; the quantum course supplied its basic principles, although even the harmonic oscillator had to be left for a later return. Light was also largely absent. That omission is now repaired, because light is intrinsically relativistic.

Most of the class has encountered special relativity before. Nevertheless, a show of hands favors rebuilding it from the beginning. That is the useful choice. The Lorentz transformation should not enter as a formula to memorize. It should emerge because operational measurements of space and time cannot satisfy both Galilean kinematics and the observed law of light propagation.

\subsection{Coordinates as physical apparatus}

Begin with a Cartesian frame: an origin, three oriented axes, and coordinates \(x,y,z\). A concrete realization would fill space with a rigid lattice of meter sticks. To identify an event rather than merely a place, each point of the lattice also carries a clock. An event is then labeled
\begin{equation}
(x,y,z,t).
\end{equation}
The clocks must be synchronized by some stated procedure; that requirement will soon become the delicate part of the construction.

\begin{definition}
A reference frame is a spatial coordinate lattice together with a distributed system of synchronized clocks. An inertial frame is a nonrotating, nonaccelerating frame in which every force-free particle moves at constant velocity.
\end{definition}

Two frames in relative uniform motion assign different coordinates to the same event. A point fixed in the moving frame has a constant coordinate \(x'\), even while its unprimed coordinate changes. In Newtonian physics the spatial relation for motion along the \(x\)-direction is
\begin{equation}
x'=x-vt.
\label{eq:sr1-galilean-space}
\end{equation}
The less conspicuous Newtonian assumption is that all well-set watches remain universally synchronized:
\begin{equation}
t'=t.
\label{eq:sr1-absolute-time}
\end{equation}

\subsection{The principle of relativity}

Relativity did not begin with Einstein. Galileo and Newton already used the principle that the laws of physics have the same form in every inertial frame. Newton's law \(F=ma\), the inverse-square law of gravity, and the Coulomb force do not acquire a new form merely because the laboratory moves uniformly.

The familiar test is a sealed railroad car. If the train moves smoothly and the windows reveal nothing outside, juggling behaves just as it does in a stationary room. Balls rise, fall, and return to the hand in the ordinary way. An experiment confined to the car cannot reveal its uniform velocity. Acceleration or rotation can be detected internally; uniform translational motion cannot.

Einstein adds one physical law to this older principle: light in vacuum propagates at
\begin{equation}
c \simeq 1.86\times 10^5\ {\rm miles/s}
   \simeq 3.00\times 10^8\ {\rm m/s}.
\end{equation}
If this is a law of nature and the laws are the same in every inertial frame, every inertial observer must measure the same \(c\).

\section{Why the Galilean Transformation Fails}

One spatial dimension is enough to expose the conflict. Draw a spacetime diagram with horizontal coordinate \(x\) and vertical coordinate \(t\). A light pulse emitted from the origin and moving to the right follows
\begin{equation}
x=ct.
\label{eq:sr1-light-right}
\end{equation}
A frame moving to the right at speed \(v\) has its origin on
\begin{equation}
x=vt,
\qquad
x'=0.
\label{eq:sr1-moving-origin}
\end{equation}

\begin{figure}[htbp]
\centering
\includegraphics[width=0.90\textwidth]{lecture_01_figure_01.png}
\caption{The right-moving light ray and the moving origin \(x=vt\), which the primed observer labels \(x'=0\).}
\lectureframe{00:15:24}
\end{figure}

Apply the Galilean rules \eqref{eq:sr1-galilean-space}--\eqref{eq:sr1-absolute-time} to the light ray:
\begin{align}
x'&=x-vt=(c-v)t,\\
t'&=t.
\end{align}
The primed observer would therefore infer
\begin{equation}
x'=(c-v)t',
\end{equation}
so right-moving light would have speed \(c-v\). For a left-moving ray \(x=-ct\),
\begin{equation}
x'=-(c+v)t',
\end{equation}
whose speed has magnitude \(c+v\).

\begin{figure}[htbp]
\centering
\includegraphics[width=0.90\textwidth]{lecture_01_figure_02.png}
\caption{The Galilean substitution predicts \(c-v\) for the right-moving light ray, contradicting frame-independent light speed.}
\lectureframe{00:17:23}
\end{figure}

Ordinary velocity addition is an excellent approximation for slow mechanical motion, but measurements of light do not show the predicted \(c-v\) and \(c+v\). The transformation law, not the light postulate, must be reconsidered. The central suspect is the apparently innocent statement \(t'=t\): it asserts an absolute notion of simultaneity.

\section{Synchronizing Distant Clocks}

To say that two clocks at different places show the same time requires an experiment. Place a third observer at their midpoint, with equal distances established by meter sticks. Instruct the two end clocks to emit flashes when each reads noon. If the flashes reach the midpoint together, call the clocks synchronized.

This is not a philosophical definition laid over the physics. It is a reproducible operation made from clocks, distances, and light signals.

Now let an observer move past the midpoint. The end flashes that arrive together at the stationary midpoint do not arrive together at the moving observer: while the light is in flight, that observer moves toward one flash and away from the other. Events judged simultaneous by the stationary clocks are therefore not simultaneous according to the moving frame's equally valid synchronization procedure.

In the unprimed frame, a line of simultaneous events is
\begin{equation}
t=\text{constant}.
\end{equation}
The corresponding primed line will be tilted in the \(x\)-\(t\) diagram. Before asking for the transformation formulas, we can construct that tilt with light signals.

\section{The Moving Frame Drawn with Light}

Set \(c=1\). One may measure distance in light-seconds when time is measured in seconds, or measure time by the light-travel time across one chosen unit of length. Either convention makes the light rays
\begin{equation}
x=\pm t,
\label{eq:sr1-light-units}
\end{equation}
and hence \(45^\circ\) lines on equally scaled axes.

Draw the light rays first. Then draw a family of observers at rest in the moving frame:
\begin{equation}
x=vt,\qquad x=vt+1,\qquad x=vt+2.
\label{eq:sr1-moving-family}
\end{equation}
Their worldlines are parallel because the observers maintain fixed separations in their own inertial lattice.

\begin{figure}[htbp]
\centering
\includegraphics[width=0.90\textwidth]{lecture_01_figure_03.png}
\caption{Light rays at \(45^\circ\) and the parallel worldlines \(x=vt\), \(x=vt+1\), and \(x=vt+2\) of the moving coordinate lattice.}
\lectureframe{00:29:55}
\end{figure}

\subsection{Fred, Mary, and Seymour}

Let Fred occupy \(x=vt\), Mary occupy \(x=vt+1\), and Seymour occupy \(x=vt+2\). At the origin \(O\), Fred sends a right-moving light signal to Mary. Seymour sends a left-moving signal so that the two flashes reach Mary at the same event \(a\). By the moving frame's midpoint synchronization rule, Seymour's emission event \(b\) is simultaneous with Fred's emission at \(O\). The line through \(O\) and \(b\) is therefore \(t'=0\).

The event \(a\) lies on both
\begin{equation}
x=t
\end{equation}
and
\begin{equation}
x=vt+1.
\end{equation}
Consequently,
\begin{equation}
t_a=x_a=\frac{1}{1-v}.
\label{eq:sr1-event-a}
\end{equation}

The left-moving light ray through \(a\) obeys
\begin{equation}
x+t=x_a+t_a=\frac{2}{1-v}.
\label{eq:sr1-ab-light}
\end{equation}
Event \(b\) also lies on Seymour's worldline,
\begin{equation}
x-vt=2.
\label{eq:sr1-seymour-line}
\end{equation}
Solving \eqref{eq:sr1-ab-light} and \eqref{eq:sr1-seymour-line} gives
\begin{align}
t_b&=\frac{2v}{1-v^2},\\
x_b&=\frac{2}{1-v^2}.
\label{eq:sr1-event-b}
\end{align}
The slope of the line from \(O\) to \(b\) is therefore
\begin{equation}
\frac{t_b}{x_b}=v,
\end{equation}
so the moving simultaneity line is
\begin{equation}
t=vx,
\qquad
t'=0.
\label{eq:sr1-moving-simultaneity}
\end{equation}

\begin{figure}[htbp]
\centering
\includegraphics[width=0.93\textwidth]{lecture_01_figure_04.png}
\caption{The Fred--Mary--Seymour synchronization construction. Equal-arrival light signals identify \(O\) and \(b\) as simultaneous in the moving frame, yielding \(t=vx\).}
\lectureframe{00:44:38}
\end{figure}

The two axes of the primed frame now have parallel forms:
\begin{align}
x'=0 &\quad\Longleftrightarrow\quad x=vt,\\
t'=0 &\quad\Longleftrightarrow\quad t=vx.
\end{align}
They are reflections of one another about the right-moving light line \(x=t\). The geometry has already replaced universal simultaneity.

\begin{classroomqa}
\audiencequestion{What happens to the construction when \(v=1\) or \(v=-1\)?}
\lecturerresponse{The denominators in \eqref{eq:sr1-event-a}--\eqref{eq:sr1-event-b} vanish, and the relevant intersections recede to infinity. At \(\lvert v\rvert=1\), a material observer would move on a light ray and could not complete the ordinary midpoint exchange. This singular behavior is the first geometric warning that \(c\) is a limiting speed. \lecturetimestamp{00:40:43}}
\end{classroomqa}

\section{Deriving the Lorentz Transformation}

The condition \(x'=0\) on \(x=vt\) requires \(x'\) to contain the factor \(x-vt\). Likewise, \(t'=0\) on \(t=vx\) requires the factor \(t-vx\). Homogeneity of spacetime makes the transformation linear, so write
\begin{align}
x'&=(x-vt)f(v),\\
t'&=(t-vx)g(v).
\label{eq:sr1-linear-ansatz}
\end{align}

\begin{figure}[htbp]
\centering
\includegraphics[width=0.90\textwidth]{lecture_01_figure_05.png}
\caption{The linear transformation is built from the two vanishing conditions, with velocity-dependent scale factors \(f(v)\) and \(g(v)\).}
\lectureframe{00:49:35}
\end{figure}

Along a right-moving light ray \(x=t\), invariance of light speed requires \(x'=t'\). Equation~\eqref{eq:sr1-linear-ansatz} then becomes
\begin{align}
x'&=(1-v)x f(v),\\
t'&=(1-v)x g(v),
\end{align}
and hence
\begin{equation}
f(v)=g(v).
\label{eq:sr1-f-equals-g}
\end{equation}
Thus one common scale factor remains:
\begin{align}
x'&=(x-vt)f(v),\\
t'&=(t-vx)f(v).
\label{eq:sr1-direct-unscaled}
\end{align}
This is essentially the economical line of argument used in Einstein's first relativity paper: operational simultaneity and invariant light speed determine the transformation.

\subsection{Reciprocity fixes the scale}

There is no distinguished \emph{really moving} inertial frame. If the primed frame moves at \(v\) relative to the unprimed frame, the unprimed frame moves at \(-v\) relative to the primed one. The inverse map must therefore have the same form with the sign of \(v\) reversed:
\begin{align}
x&=(x'+vt')f(v),\\
t&=(t'+vx')f(v).
\label{eq:sr1-inverse-unscaled}
\end{align}

\begin{figure}[htbp]
\centering
\includegraphics[width=0.90\textwidth]{lecture_01_figure_06.png}
\caption{Reciprocity gives the inverse transformation by reversing the relative velocity while retaining the same scale factor.}
\lectureframe{00:55:12}
\end{figure}

Insert \eqref{eq:sr1-direct-unscaled} into the first of \eqref{eq:sr1-inverse-unscaled}:
\begin{align}
x
&=\left[(x-vt)f(v)+v(t-vx)f(v)\right]f(v)\\
&=x(1-v^2)f(v)^2.
\end{align}
Because this composition must be the identity,
\begin{equation}
(1-v^2)f(v)^2=1,
\end{equation}
and therefore
\begin{equation}
f(v)=\pm\frac{1}{\sqrt{1-v^2}}.
\label{eq:sr1-factor-roots}
\end{equation}

\begin{figure}[htbp]
\centering
\includegraphics[width=0.90\textwidth]{lecture_01_figure_07.png}
\caption{Composing the direct and inverse maps fixes the magnitude of the common factor \(f(v)\).}
\lectureframe{00:59:31}
\end{figure}

\begin{classroomqa}
\audiencequestion{Why discard the negative square root in \eqref{eq:sr1-factor-roots}?}
\lecturerresponse{At zero relative velocity the transformation must continuously approach the identity, \(x'=x\) and \(t'=t\). The negative branch instead gives \(x'=-x\) and \(t'=-t\), a reversal of both coordinate orientations. Choosing it would introduce an abrupt reflection rather than a continuously moving frame. \lecturetimestamp{00:59:27}}
\end{classroomqa}

With the positive branch, the Lorentz transformation in units \(c=1\) is
\begin{align}
x'&=\frac{x-vt}{\sqrt{1-v^2}},&
t'&=\frac{t-vx}{\sqrt{1-v^2}}.
\label{eq:sr1-lorentz-c1}
\end{align}

\begin{figure}[htbp]
\centering
\includegraphics[width=0.90\textwidth]{lecture_01_figure_08.png}
\caption{The completed Lorentz transformation in units where the speed of light is one.}
\lectureframe{01:02:05}
\end{figure}

\subsection{Restoring Ordinary Units}

Dimensional analysis restores ordinary units. The denominator can depend only on the dimensionless ratio \(v^2/c^2\), while \(vx\) in the time equation must become \(vx/c^2\):
\begin{align}
x'&=\gamma(x-vt),&
t'&=\gamma\left(t-\frac{vx}{c^2}\right),
\label{eq:sr1-lorentz-dimensional}\\
\gamma&=\frac{1}{\sqrt{1-v^2/c^2}}.
\label{eq:sr1-gamma}
\end{align}
The inverse transformation is
\begin{align}
x&=\gamma(x'+vt'),&
t&=\gamma\left(t'+\frac{vx'}{c^2}\right).
\label{eq:sr1-lorentz-inverse-dimensional}
\end{align}

\begin{figure}[htbp]
\centering
\includegraphics[width=0.90\textwidth]{lecture_01_figure_09.png}
\caption{The Lorentz formulas with \(c\) restored: \(v^2/c^2\) controls the scale and \(vx/c^2\) mixes position into time.}
\lectureframe{01:05:01}
\end{figure}

For \(v\ll c\),
\begin{equation}
\gamma=1+\frac{1}{2}\frac{v^2}{c^2}+O\!\left(\frac{v^4}{c^4}\right),
\end{equation}
and the transformation reduces to
\begin{equation}
x'\approx x-vt,
\qquad
t'\approx t.
\end{equation}
At terrestrial speeds the ratios are usually tiny, explaining the success of Newtonian kinematics. Near \(c\), the corrections become large.

The motion was chosen along the \(x\)-axis. A transverse measuring arm cannot distinguish one inertial frame from the other without violating reciprocity, so
\begin{equation}
y'=y,
\qquad
z'=z.
\end{equation}

\section{Measuring a Moving Rod}

Length means the separation of the two endpoints measured at the same time in the frame doing the measurement. It does not mean the apparent angular size of a photographed object, because light-travel delays introduce a separate optical problem.

Let a rod be at rest in the unprimed frame with endpoints
\begin{equation}
x=0,
\qquad
x=1.
\end{equation}
The primed observer measures the endpoint events on a single primed simultaneity slice, say \(t'=0\). At the far endpoint \(x=1\),
\begin{equation}
t'=0
\quad\Longrightarrow\quad
t=\frac{vx}{c^2}.
\end{equation}
Temporarily returning to \(c=1\), this event has \(t=v\), and
\begin{equation}
x'=\frac{x-vt}{\sqrt{1-v^2}}
   =\frac{1-v^2}{\sqrt{1-v^2}}
   =\sqrt{1-v^2}.
\end{equation}

\begin{figure}[htbp]
\centering
\includegraphics[width=0.90\textwidth]{lecture_01_figure_10.png}
\caption{The worldlines \(x=0\) and \(x=1\) of a rod at rest in the unprimed frame; a moving observer must compare them on one \(t'\)-slice.}
\lectureframe{01:11:14}
\end{figure}

\begin{figure}[htbp]
\centering
\includegraphics[width=0.90\textwidth]{lecture_01_figure_11.png}
\caption{The simultaneous primed-frame endpoint measurement gives the contraction factor \(\sqrt{1-v^2}\) in units \(c=1\).}
\lectureframe{01:16:03}
\end{figure}

Restoring units and a general rest length \(L_0\),
\begin{equation}
L=\frac{L_0}{\gamma}
 =L_0\sqrt{1-\frac{v^2}{c^2}}.
\label{eq:sr1-length-contraction}
\end{equation}

\begin{classroomqa}
\audiencequestion{If the far end of the rod is at \(x=1\), should its coordinate not also be \(x'=1\)?}
\lecturerresponse{No. The rod is at rest in the unprimed frame, so \(x=0\) and \(x=1\) label its fixed endpoints there. The primed coordinate of the far endpoint changes with time. Its value on the measuring slice \(t'=0\) is precisely what the Lorentz transformation computes. \lecturetimestamp{01:14:24}}
\end{classroomqa}

\begin{classroomqa}
\audiencequestion{How can each inertial observer consistently say that the other observer's rod is shorter?}
\lecturerresponse{They measure different pairs of endpoint events. The unprimed observer requires equal \(t\); the primed observer requires equal \(t'\). Because simultaneity is frame-dependent, those two endpoint pairs are not the same pair of events. Reciprocity therefore produces no contradiction. \lecturetimestamp{01:16:28}}
\end{classroomqa}

The operational version is worth making explicit. Place synchronized clocks and assistants along the measuring frame. At a prescribed local time, each assistant records the position of the endpoint passing nearby. Subtracting those two simultaneously recorded positions gives the rod's length in that frame. Nothing depends on how the rod happens to look.

\begin{classroomqa}
\audiencequestion{When \(c=1\), does a unit rod have to be one light-second long?}
\lecturerresponse{Only if time is measured in seconds. One may instead use meters for length and measure time in the light-travel time across one meter. Setting \(c=1\) ties the units of time and length together; it does not prescribe which physical unit must be chosen first. \lecturetimestamp{01:21:17}}
\end{classroomqa}

\section{A Moving Clock and the Twin Question}

Now follow a clock at rest in the primed frame. Along its worldline,
\begin{equation}
x'=0.
\end{equation}
Compare the event at which it reads \(t'=1\) with the origin. The inverse transformation gives
\begin{equation}
t=\gamma\left(t'+\frac{vx'}{c^2}\right)=\gamma.
\end{equation}
Thus one unit of elapsed time on the moving clock corresponds to \(\gamma\) units of unprimed coordinate time:
\begin{equation}
\Delta t=\gamma\,\Delta\tau,
\qquad
\Delta\tau=\frac{\Delta t}{\gamma}.
\label{eq:sr1-time-dilation}
\end{equation}

\begin{figure}[htbp]
\centering
\includegraphics[width=0.90\textwidth]{lecture_01_figure_12.png}
\caption{A clock carried on \(x'=0\) records less elapsed time between the same two events than the stationary coordinate time: \(t=\gamma t'\).}
\lectureframe{01:26:42}
\end{figure}

The effect concerns every physical clock, not merely a mechanism labeled as one. A wristwatch, an atomic transition, metabolism, and biological aging all follow the same proper time. A cheap watch bought on a New York street would not escape relativity by being a poor Rolex; it might keep inaccurate time, but whatever process it uses is still subject to the same kinematics.

\begin{classroomqa}
\audiencequestion{If each inertial observer sees the other's clock running slowly, why does the traveling twin return younger?}
\lecturerresponse{A round trip is not a comparison of two globally equivalent inertial worldlines. The traveler must turn around and join a different inertial frame, while the stay-at-home twin remains on one inertial worldline. The turnaround breaks the exchange symmetry, so the accumulated proper times can differ. \lecturetimestamp{01:26:37}}
\end{classroomqa}

\begin{classroomqa}
\audiencequestion{Is the acceleration during the turnaround itself what causes the age difference?}
\lecturerresponse{Acceleration marks where the symmetry is broken, but the age difference is calculated by integrating proper time along the complete paths. The turnaround may be made brief or treated as a change between inertial legs; the unequal spacetime paths still have unequal proper times. \lecturetimestamp{01:29:20}}
\end{classroomqa}

\section{The Invariant Geometry}

The time-dilation diagram looks perverse if it is read with Euclidean geometry. The slanted path seems like a hypotenuse, yet it carries less proper time than the vertical path. The resolution is not a trick in the algebra. Spacetime has a different metric.

In a Euclidean plane, a rotation changes \(x\) and \(y\) but preserves
\begin{equation}
x'^2+y'^2=x^2+y^2.
\label{eq:sr1-euclidean-invariant}
\end{equation}
It is natural to test whether a Lorentz transformation preserves \(t^2+x^2\). It does not. In units \(c=1\),
\begin{align}
t'^2
&=\frac{t^2-2vtx+v^2x^2}{1-v^2},\\
x'^2
&=\frac{x^2-2vtx+v^2t^2}{1-v^2}.
\end{align}
Adding leaves the wrong coefficients. Subtracting instead cancels the mixed terms and gives
\begin{align}
t'^2-x'^2
&=
\frac{t^2-2vtx+v^2x^2-x^2+2vtx-v^2t^2}{1-v^2}\\
&=t^2-x^2.
\label{eq:sr1-interval-invariance}
\end{align}

\begin{figure}[htbp]
\centering
\includegraphics[width=0.90\textwidth]{lecture_01_figure_13.png}
\caption{Subtracting the squared transformed coordinates cancels the cross terms and exposes the invariant \(t^2-x^2\).}
\lectureframe{01:38:07}
\end{figure}

Coordinates are frame-dependent components. The combination
\begin{equation}
\Delta s^2=c^2\Delta t^2-\Delta x^2
\label{eq:sr1-spacetime-interval}
\end{equation}
is the same in every inertial frame. With more spatial dimensions,
\begin{equation}
\Delta s^2
=c^2\Delta t^2-\Delta x^2-\Delta y^2-\Delta z^2.
\end{equation}

\begin{classroomqa}
\audiencequestion{Why is it legitimate to abandon the ordinary geometric rule that a hypotenuse is longer?}
\lecturerresponse{That rule belongs to the positive-definite Euclidean metric. The coordinates here were defined operationally with rods, clocks, and light signals, and their transformation preserves a difference of squares. Minkowski geometry is not assumed to evade the paradox; it is discovered from those measurements. \lecturetimestamp{01:43:21}}
\end{classroomqa}

\subsection{Proper time}

For two events on the worldline of a clock, choose the clock's rest frame. There \(\Delta x'=0\), so invariance gives
\begin{equation}
c^2\Delta\tau^2
=c^2\Delta t^2-\Delta x^2,
\end{equation}
or
\begin{equation}
\Delta\tau
=\sqrt{\Delta t^2-\frac{\Delta x^2}{c^2}}.
\label{eq:sr1-proper-time}
\end{equation}
This is the actual time recorded by the clock between the two events.

\begin{figure}[htbp]
\centering
\includegraphics[width=0.90\textwidth]{lecture_01_figure_14.png}
\caption{In the clock's rest frame the spatial separation vanishes, so the invariant interval becomes the square of its proper time.}
\lectureframe{01:41:28}
\end{figure}

\begin{classroomqa}
\audiencequestion{Should proper time be \(t^2-x^2/c^2\), or its square root?}
\lecturerresponse{The difference of squares is the square of the proper time. The elapsed time read by a clock is \(\tau=\sqrt{t^2-x^2/c^2}\) for separations measured from a common origin. The missing square root in the initial wording must be restored. \lecturetimestamp{01:51:02}}
\end{classroomqa}

\subsection{Timelike, null, and spacelike separation}

Two distinct events can have zero interval. In units \(c=1\),
\begin{equation}
\Delta t^2-\Delta x^2=0
\quad\Longleftrightarrow\quad
\Delta x=\pm\Delta t.
\end{equation}
Such events are connected by a light ray.

\begin{figure}[htbp]
\centering
\includegraphics[width=0.90\textwidth]{lecture_01_figure_15.png}
\caption{Null-separated events lie on \(x=\pm t\), where the spacetime interval and the proper time of the lightlike path vanish.}
\lectureframe{01:46:40}
\end{figure}

If light could carry a conventional clock, no proper-time tick would occur between emission and absorption. This is not a valid rest frame for light; it is the limiting statement that every null segment has zero proper time.

The light rays divide spacetime into three invariant classes:
\begin{align}
c^2\Delta t^2-\Delta x^2&>0
&&\text{timelike},\\
c^2\Delta t^2-\Delta x^2&=0
&&\text{null},\\
c^2\Delta t^2-\Delta x^2&<0
&&\text{spacelike}.
\end{align}
A physical clock can connect timelike-separated events, and the positive interval determines its proper time. No subluminal object can connect spacelike-separated events. For a spacelike pair one may choose a frame in which the events are simultaneous and define the proper distance
\begin{equation}
\ell_{\rm proper}
=\sqrt{\Delta x^2-c^2\Delta t^2}.
\end{equation}

\begin{figure}[htbp]
\centering
\includegraphics[width=0.90\textwidth]{lecture_01_figure_16.png}
\caption{The light cone separates timelike pairs that a clock can connect from spacelike pairs that no subluminal object can connect.}
\lectureframe{01:48:47}
\end{figure}

\begin{classroomqa}
\audiencequestion{When should the invariant be called proper time, and when should it be called proper length?}
\lecturerresponse{For a timelike pair there is a rest frame in which both events occur at one place, and the interval gives the proper time. For a spacelike pair there is a frame in which both occur at one time, and the opposite-sign invariant gives their proper distance. Null separation is the boundary between the two. \lecturetimestamp{01:47:22}}
\end{classroomqa}

\begin{classroomqa}
\audiencequestion{Where does \(c\) belong when the units are not chosen so that \(c=1\)?}
\lecturerresponse{The invariant is \(c^2\Delta t^2-\Delta x^2\), with dimensions of length squared. Dividing by \(c^2\) gives \(\Delta\tau^2=\Delta t^2-\Delta x^2/c^2\), with dimensions of time squared. \lecturetimestamp{01:49:56}}
\end{classroomqa}

\section{Lorentz Symmetry and Its Historical Meaning}

\begin{classroomqa}
\audiencequestion{Can a Lorentz transformation be regarded as a rotation?}
\lecturerresponse{Yes, but not as an ordinary Euclidean rotation. It is a hyperbolic rotation preserving \(t^2-x^2\). Historically one can describe it as a rotation through an imaginary angle; rapidity gives the modern real parametrization. \lecturetimestamp{01:51:40}}
\end{classroomqa}

Einstein's route was tied closely to electromagnetism. According to his later recollection, already as a teenager he wondered what would happen if he chased a light wave. Galilean velocity subtraction would make the wave appear stationary, but Maxwell's equations have no corresponding frozen oscillatory light wave in vacuum. That conflict pointed to the kinematics, not to a peculiar mechanical property of measuring rods.

\begin{classroomqa}
\audiencequestion{Was Einstein's argument primarily a response to experiments such as Michelson--Morley?}
\lecturerresponse{The experimental background mattered historically, but Einstein's own testimony emphasized the structure of Maxwell's equations and the thought experiment of chasing a light beam. How much he knew or relied on particular interferometer results is historically debated, so the clean physical point is the symmetry conflict between Maxwell and Galilean transformations. \lecturetimestamp{01:52:04}}
\end{classroomqa}

Maxwell's equations are invariant under Lorentz transformations rather than Galilean ones. Lorentz and FitzGerald had already found related contraction formulas while retaining an ether and interpreting contraction as a dynamical deformation caused by motion through it. On that view the effect might have been an approximate material response with further corrections.

Einstein's decisive step was to make the relativity principle and invariant light speed exact kinematic principles. Length contraction and time dilation then apply to every suitable rod and clock because they express the geometry relating inertial frames, not a special force exerted by an ether. The next task is to formulate particle motion inside this geometry.
